% Chapter 4: Phase One
% Scene function: First work night. Tert maps the house's acoustic
% skeleton and makes one small environmental change. The craft in
% action. The tedium of excellence.

\chapter{Fourteen Minutes}

The subject goes to bed at eleven-fourteen.

I know this because I have been in the house since six, positioned in the hallway between the bathroom and the sickroom door, and I have been tracking his movements by sound for five hours. This is not exciting work. It is necessary work. Phase 1 requires a complete baseline of the subject's evening routine: when he eats, when he watches television, when he moves from room to room, when the lights go off. You need the baseline before you can introduce variation, the same way a musician needs to hear the instrument in tune before they can play it out of tune on purpose.

He ate dinner at seven. Something involving a pot and a cutting board and the sound of a single plate being set on a table. He watched television from seven-forty until nine-thirty. I could hear it through the wall but could not tell what it was. He washed the pot and the plate and the cutting board. He brushed his teeth. He walked past the sickroom door without pausing, which is what he did the first time I was here and what I expect he does every time. He turned off the hallway light. He went to bed. The bedroom light stayed on for another twenty minutes and then went off.

Eleven-fourteen.

I wait another forty minutes. You do not begin until the subject is in deep sleep. This is not a guess. Breathing patterns change. The house changes around a sleeping person: it becomes the only thing making sound. A sleeping human generates a rhythm, and after forty minutes that rhythm has stabilized. The house is now the subject's instrument. My job is to play it.

I start in the basement.

The basement is where the joists are exposed, which means I can see the structural skeleton of the house. Every creak the subject hears on the floor above originates from one of these joists expanding or contracting against its bearing point. The house has a pattern. It creaks at predictable intervals based on the temperature curve through the night: the wood contracts as the exterior cools, the contraction generates stress at the joints, and the stress releases as a sound the subject has been hearing every night for twenty years.

I am going to make that sound happen six minutes early.

Not tonight. Tonight I am mapping. I am lying on the concrete floor of the basement with my hand on a joist, feeling the thermal contraction through the wood, counting the interval between stress events. This is the part of the job nobody talks about in briefings, because it is not dramatic and it does not photograph well for the departmental newsletter. There is no departmental newsletter. But if there were, this would not make the cover.

The interval between creaks on this joist is roughly fourteen minutes. The next joist over is closer to nine. The one nearest the living room, directly under the bookshelf wall, has an irregular pattern that I will need two more nights to map.

I spend four hours in the basement. The concrete is not comfortable. I have never filed the form for this, and I suspect nobody has.

At three in the morning I come upstairs. The subject is still asleep. I can hear him breathing through the bedroom door, which he does not close all the way. People who have been caregivers often sleep with doors open. It is a habit that outlasts its reason.

I stand in the living room in the dark. The room is different at night. The books on the shelves are dark shapes. The fireplace is a black rectangle. The two armchairs are silhouettes, one facing slightly away from the other. The empty nail above the mantel catches the light from the street lamp outside in a way that makes it the brightest point in the room. I did not plan this. The house did it on its own.

I do one thing before I leave. A small thing. I have been in the house for nine hours without introducing any change to the environment, which is correct for a first work night. But I allow myself one small adjustment, something so minor it barely qualifies.

I shift the thermostat down two degrees.

Not enough for the subject to wake up cold. Not enough for him to notice in the morning and check the setting. Just enough that when his body moves through the hallway at six a.m. on the way to the bathroom, the air will be slightly different from yesterday, and the difference will register below the threshold of conscious thought as a sensation he cannot name and will not remember having.

This is what Phase 1 is. Not fear. Not dread. A whisper of wrongness so faint that the subject's only response is to feel, for a moment, slightly less at home in the place where he lives.

I let myself out the back door. The sky is turning gray. A cat on the porch across the street watches me without interest, which is typical. Animals do not react to us the way the movies suggest. They notice us the same way they notice a change in the weather: briefly, and then they go back to what they were doing.
